\subsection{固定楕円パラメータにおける全変形の中心部分環}
\label{subsec:cubic-third-root-unit-boundary-L-polynomial-center}
\begin{theorem}[$E=2$、全 $u,v$ の中心自由加群]
\label{thm:cubic-third-root-unit-boundary-L-polynomial-center}
六次代数体 $K$ と $E=2$ を固定する。
$u,v$ を独立変数として、37個の四生成元の語からなる
多項式係数の中心元 $z_9(u,v)$ を選べる。
$k$ を標数0代数閉体とすると、全六埋込みと全 $u,v\in k$ について
\[
 k[x,\Psi,z_9(u,v)]\subset Z(A_{u,v})
\]
は次数 $1,3,9$ の多項式環であり、$A_{u,v}$ はその上自由階数81である。
さらに $\dim_k Z(A_{u,v})_9=5$ が全 $u,v$ で成り立つ。
この中心部分環が全中心であること、全点での素性・PI次数9・H3排除は、
本定理には含めない。
\end{theorem}
\begin{proof}
係数環を $P=K[u,v]$ とする。
元の全二次関係と $G$ の挿入から、$A_9,A_{10}$ の語基底を逐次構成する。
$x$ の出現回数を先に比較する順序を使うと、全relation pivotは
$u,v$ に依存しない非零 $K$ 元となる。
従って各次数の商は $P$ 上自由で、全ての特殊化と両側乗法を正確に記述する。
次数9,10の階数は136,166である。

四生成元との交換子を一つの行列にまとめる。
ここでも $u,v$ に依存しない単元成分だけをpivotとして選ぶ消去が
131回進み、残る全行は多項式として零となる。
従って中心核は $P$ 上自由階数5で、全ての特殊化後も次元5である。
五つの核ベクトルの四交換子を元の左右乗法へ再代入して零を確認する。

$x^9,x^6\Psi,x^3\Psi^2,\Psi^3$ に対して、
$u=1,v=0$ で独立になる核ベクトルを一つ選び、$z_9$ とする。
その $x=0$ の像には $u,v$ が全く現れないことも直接確認する。
従って $x=\Psi=0$ での $z_9$ の像は、固定例の既証の非零正則元の
非零scalar倍であり、全 $u,v$ に共通である。
実際、新しい $z_9$ を用いて境界商を元から計算しても、
三つの中心元によるArtinian商の次数0から13の次元は
\[
 1,3,6,8,9,9,9,9,9,8,6,3,1,0
\]
となる。

$x,\Psi$ の正則性は全変形で既証であり、
残る $z_9$ の正則性、中心三元の代数的独立性、自由階数81は
固定例と同じHilbert級数と次数付きNakayamaの議論で従う。
ここで主張するのは全ての幾何学的パラメータ点についての結論であり、
全中心の任意底変更やintrinsic moduliの降下まで主張していない。

再現器 \texttt{verify\_cubic\_third\_root\_unit\_boundary\_L\_polynomial\_center.py}
は元の関係を左右から全基底に作用させ、パラメータ多項式のまま次数10まで、
Artinian商で次数13まで零を確認する。
JSONに $z_9$ の全語係数を保存する。
変形パラメータの非定数多項式を割る操作は行っていない。
\end{proof}
